% !PRINTED 622
\weberpage{622}
\begin{center}
\S.\,171.

\gesperrt{The Prime Ideals in the Field $\Omega_m$.}
\end{center}

% [sic: The printed reference 164 should be 169, theorem I and (7), where the prime-power factorisation is established.]
% [Source shorthand: q is associated with sigma^mu, not generally equal to it as an element; the corresponding principal ideals are equal.]
We have seen in \S.\,164 that the natural prime $q$
is the $\mu$th power of a prime $\sigma$ in the field $\Omega_m$.

In order to determine all the prime ideals that exist in the
field $\Omega_m$, it therefore remains to decompose into prime factors
all the natural primes $p$ different from $q$, and to select from
among them those that are not associated.

The foundation for this investigation is formed by the theorems
of \S.\,162, when the field there denoted by $R$ is replaced by
the field of rational numbers.

Every number $\omega$ of the field $\Omega_m$ is carried by one of the $\mu$ substitutions
\[
s_n = (r, r^n)
\]
into a determinate other number $\omega_n$, which is likewise contained
in $\Omega_m$. If
\weq{1}{\omega = \omega_1 = a_0 + a_1 r + a_2 r^2 + \cdots + a_{\mu-1} r^{\mu-1},}
then
\weq{2}{\omega_n = a_0 + a_1 r^n + a_2 r^{2n} + \cdots + a_{\mu-1} r^{(\mu-1)n}.}
If one of these numbers $\omega_n$ is an integer, as we shall now
assume, then so are all the others, and this occurs always then
and only then, when the rational numbers
$a_0, a_1, \ldots, a_{\mu-1}$ are integers.

If $h, k$ are any two exponents, then
\[
r^h - r^k = r^k(r^{h-k} - 1),
\]
and by \S.\,169, (9) this is associated with a power of $\sigma$,
hence, if $\mathfrak{p}$ denotes a prime ideal dividing $p$, it is not
divisible by $\mathfrak{p}$, except when $h \equiv k \pmod{m}$. From this it results
% [sic: The following assertion needs congruence for every integral omega, not for one fixed omega (omega=1 is a counterexample). Choosing omega=r proves the required assertion about the inertia group X, so the unramifiedness conclusion remains valid.]
that two numbers $\omega_h, \omega_k$ can be congruent with respect to the modulus $\mathfrak{p}$ only when $h \equiv k \pmod{m}$, and that therefore the
group $X$ of \S.\,162, which satisfies the condition
\[
\omega\,|\,\chi \equiv \omega \pmod{\mathfrak{p}}
\]
contains only the \gesperrt{identical substitution}. But from this it
follows that $g = 1$, i.\,e. that $p$ \gesperrt{is not divisible by the square
of a prime ideal.}
